\documentclass[11pt,a4paper,oneside]{report}

% ==============================================================================
% CHARGEMENT DES PACKAGES DE STYLE ET MACROS (COMPATIBLE 100% OVERLEAF)
% ==============================================================================
\usepackage{fascicule_style}
\usepackage{macros}

% Métadonnées du fascicule
\setfiliere{3\textsuperscript{ème} Prépa-Métiers}
\setmatiere{Mathématiques}
\setannee{2025--2026}
\setlycee{Lycée Professionnel \& Métiers}
\setprof{Équipe Pédagogique Maths-Sciences}

\begin{document}

% ==============================================================================
% PAGE DE COUVERTURE MODERNE & ÉLÉGANTE (DESIGN TIKZ)
% ==============================================================================
\begin{titlepage}
\begin{tikzpicture}[remember picture, overlay]
    % Fond et dégradé / bandes graphiques
    \fill[navy900] (current page.north west) rectangle (current page.south east);
    \fill[colPrimary] (current page.north west) rectangle ([xshift=1.2cm]current page.south west);
    \fill[colSecondary!80!navy900] ([yshift=-8cm]current page.north west) -- ([xshift=14cm, yshift=-8cm]current page.north west) -- ([xshift=6cm, yshift=-16cm]current page.north west) -- ([yshift=-16cm]current page.north west) -- cycle;
    
    % Encart supérieur
    \node[anchor=north west, text=slate300, font=\sffamily\bfseries\large] at ([xshift=2.2cm, yshift=-2.2cm]current page.north west) {
        \thelycee\ -- Année Scolaire \theannee
    };
    
    % Badge Filière
    \node[anchor=north west] at ([xshift=2.2cm, yshift=-3.5cm]current page.north west) {
        \tcbox[
            colback=colSecondary,
            colframe=colSecondary,
            arc=6pt,
            boxrule=0pt,
            top=4pt, bottom=4pt, left=12pt, right=12pt,
            fontupper=\sffamily\bfseries\color{navy900}\large
        ]{RÉFÉRENTIEL NATIONAL $\cdot$ DNB PRO}
    };
    
    % Titre Principal
    \node[anchor=north west, text=white, font=\sffamily\bfseries\fontsize{38}{44}\selectfont] at ([xshift=2.2cm, yshift=-5.2cm]current page.north west) {
        FASCICULE DE COURS\\\& EXERCICES
    };
    
    % Sous-Titre Matière & Classe
    \node[anchor=north west, text=colSecondary, font=\sffamily\bfseries\fontsize{24}{28}\selectfont] at ([xshift=2.2cm, yshift=-8.5cm]current page.north west) {
        \thematiere\ -- \thefiliere
    };
    
    % Descriptif du contenu
    \node[anchor=north west, text=slate200, font=\sffamily\normalsize, text width=14cm, align=left] at ([xshift=2.2cm, yshift=-10.2cm]current page.north west) {
        \textbf{Le livret complet pour réussir son année et l'épreuve du Brevet :}
        \begin{itemize}
            \item \textbf{Cours structuré} avec définitions, propriétés et méthodes modèles
            \item \textbf{Questions Flash \& Automatismes} de calcul mental
            \item \textbf{Travaux Dirigés (TD)} progressifs et problèmes contextualisés
            \item \textbf{TP TICE} (Tableur, GeoGebra, Scratch)
            \item \textbf{Annales officielles DNB Pro} avec corrigés pas à pas
        \end{itemize}
    };
    
    % Cartouche Auteur / Nom de l'élève
    \node[anchor=south west] at ([xshift=2.2cm, yshift=2.2cm]current page.south west) {
        \begin{tcolorbox}[
            enhanced,
            width=16.5cm,
            colback=slate800,
            colframe=colSecondary,
            boxrule=1.2pt,
            arc=6pt,
            fonttitle=\sffamily\bfseries\small\color{colSecondary},
            title={INFORMATIONS ÉLÈVE},
            coltext=white
        ]
        \sffamily
        \textbf{Nom \& Prénom :} \pointille[8cm] \hfill \textbf{Classe :} 3\textsuperscript{ème} PM\dots \\[2mm]
        \textbf{Professeur :} \theprof \hfill \textbf{Groupe :} \pointille[3cm]
        \end{tcolorbox}
    };
\end{tikzpicture}
\end{titlepage}

% ==============================================================================
% SOMMAIRE / TABLE DES MATIÈRES INTERACTIVE
% ==============================================================================
\pagenumbering{roman}
\tableofcontents
\newpage

\pagenumbering{arabic}
\setcounter{page}{1}

% ==============================================================================
% INCLUSION MODULAIRE DES CHAPITRES
% ==============================================================================
\chapter*{Progression Annuelle \& Référentiel}
\addcontentsline{toc}{chapter}{Progression Annuelle \& Référentiel}
\markboth{PROGRESSION ANNUELLE}{RÉFÉRENTIEL 3PM}

\begin{cadreretenu}[Objectifs de la Classe de 3\textsuperscript{ème} Prépa-Métiers]
La classe de 3\textsuperscript{ème} Prépa-Métiers permet aux élèves de consolider le socle commun de connaissances, de compétences et de culture tout en découvrant des environnements professionnels variés. En mathématiques, l'accent est mis sur :
\begin{itemize}
    \item La consolidation des automatismes de calcul et du raisonnement logique.
    \item La résolution de situations concrètes issues des métiers (BTP, Industrie, Tertiaire, Artisanat).
    \item L'utilisation efficace des outils numériques (Calculatrice, Tableur, GeoGebra, Scratch).
    \item La préparation rigoureuse à l'épreuve de Mathématiques du \textbf{Diplôme National du Brevet (DNB Pro)}.
\end{itemize}
\end{cadreretenu}

\vspace{0.5cm}

\section*{Grille Nationale des Compétences Évaluées}

\noindent
\begin{tabularx}{\linewidth}{|c|l|X|}
\hline
\rowcolor{slate800} \textcolor{white}{\textbf{Code}} & \textcolor{white}{\textbf{Compétence}} & \textcolor{white}{\textbf{Description synthétique}} \\
\hline
\rowcolor{slate50} \capSApproprier & \textbf{S'approprier} & Extraire l'information utile, reformuler une situation, identifier les données d'un problème métier. \\
\hline
\rowcolor{white} \capAnalyser & \textbf{Analyser / Raisonner} & Choisir un modèle (linéaire, géométrique, statistique), émettre une conjecture, élaborer une démarche. \\
\hline
\rowcolor{slate50} \capRealiser & \textbf{Réalisér} & Effectuer un calcul numérique ou algébrique, utiliser la calculatrice, construire une figure, exécuter un algorithme. \\
\hline
\rowcolor{white} \capValider & \textbf{Valider} & Contrôler la vraisemblance d'un résultat, vérifier des unités de mesure, critiquer une hypothèse. \\
\hline
\rowcolor{slate50} \capCommuniquer & \textbf{Communiquer} & Rédiger une solution claire, argumenter à l'oral ou à l'écrit, expliciter les étapes de son raisonnement. \\
\hline
\end{tabularx}

\vspace{0.8cm}

\section*{Tableau de Progression Annuelle (5 Périodes)}

\noindent
\begin{tabularx}{\linewidth}{|c|l|X|c|}
\hline
\rowcolor{colPrimary} \textcolor{white}{\textbf{Période}} & \textcolor{white}{\textbf{Domaine}} & \textcolor{white}{\textbf{Chapitres \& Notions au Programme}} & \textcolor{white}{\textbf{Éval.}} \\
\hline
\rowcolor{colDefBg} \textbf{Période 1} & \textbf{Nombres \& Calculs} & $\bullet$ Divisibilité, décomposition en facteurs premiers, fractions irréductibles.\newline $\bullet$ Puissances de 10, écriture scientifique et ordres de grandeur.\newline $\bullet$ \textbf{Calcul Littéral} : conventions, développement $k(a+b)$, $(a+b)(c+d)$, réduction. & DS 1 \\
\hline
\rowcolor{white} \textbf{Période 2} & \textbf{Algèbre \& Géométrie} & $\bullet$ \textbf{Équations du 1\textsuperscript{er} degré} : résolution algébrique $ax+b=c$, problèmes concrets.\newline $\bullet$ Proportionnalité, pourcentages d'évolution, ratios et échelles.\newline $\bullet$ \textbf{Théorème de Pythagore} : calcul de longueurs, réciproque et contraposée. & DS 2 \\
\hline
\rowcolor{colDefBg} \textbf{Période 3} & \textbf{Fonctions \& Espace} & $\bullet$ \textbf{Fonctions affines et linéaires} : $f(x)=ax+b$, représentations graphiques, lecture de coefficients.\newline $\bullet$ \textbf{Trigonométrie} dans le triangle rectangle ($\cos, \sin, \tan$).\newline $\bullet$ Repérage dans le plan et dans l'espace (coordonnées 2D/3D). & Brevet Blanc 1 \\
\hline
\rowcolor{white} \textbf{Période 4} & \textbf{Géométrie \& Données} & $\bullet$ \textbf{Théorème de Thalès} \& Agrandissement-Réduction (rapports $k, k^2, k^3$).\newline $\bullet$ \textbf{Statistiques à 1 variable} : moyenne pondérée, médiane, étendue, fréquences.\newline $\bullet$ \textbf{Probabilités} : expériences aléatoires simples et à 2 épreuves, arbres. & DS 3 \\
\hline
\rowcolor{colDefBg} \textbf{Période 5} & \textbf{Grandeurs \& Synthèse} & $\bullet$ \textbf{Aires, Volumes \& Conversions} : prismes, cylindres, cônes, pyramides, sphères.\newline $\bullet$ \textbf{Algorithmique \& Scratch} : boucles, conditions, variables, déplacements.\newline $\bullet$ \textbf{Entraînement intensif aux annales du DNB Pro}. & DNB PRO \\
\hline
\end{tabularx}

\newpage


\chapter{Calcul Littéral \& Expressions}
\label{chap:calcul_litteral}

% ==============================================================================
% 1. OBJECTIFS & COMPÉTENCES DU RÉFÉRENTIEL
% ==============================================================================
\begin{cadreretenu}[Objectifs \& Compétences du DNB Pro]
\begin{itemize}
    \item \textbf{Connaissances} : Conventions d'écriture littérale, vocabulaire (somme, produit, terme, facteur), identité remarquable $a^2 - b^2 = (a-b)(a+b)$.
    \item \textbf{Capacités} :
    \begin{itemize}
        \item \capRealiser\ Réduire une expression algébrique en regroupant les termes de même nature.
        \item \capRealiser\ Développer une expression en utilisant la distributivité simple $k(a+b)$ et double $(a+b)(c+d)$.
        \item \capRealiser\ Factoriser une expression simple (facteur commun évident ou différence de deux carrés).
        \item \capAnalyser\ Traduire une situation concrète (périmètre, aire, coût de production) par une formule littérale.
        \item \capValider\ Tester l'égalité de deux expressions en substituant des valeurs numériques.
    \end{itemize}
\end{itemize}
\end{cadreretenu}

\vspace{0.3cm}

% ==============================================================================
% 2. ACTIVITÉ D'INVESTIGATION
% ==============================================================================
\section{Activité d'Introduction : L'Atelier de Menuiserie}

\begin{activite}{Conception d'une structure en bois modulable}
Un menuisier conçoit des cadres en bois rectangulaires dont la largeur est notée $x$ (en mètres) et la longueur mesure $3\text{ m}$ de plus que la largeur.

\begin{center}
\begin{tikzpicture}[scale=0.9]
    \draw[fill=amber!20, draw=amber!80!black, line width=1.2pt] (0,0) rectangle (6,3.5);
    \draw[<->, >=stealth, slate700, thick] (0,-0.4) -- (6,-0.4) node[midway, below] {Longueur : $L = x + 3$};
    \draw[<->, >=stealth, slate700, thick] (-0.4,0) -- (-0.4,3.5) node[midway, left] {Largeur : $\ell = x$};
    \node at (3,1.75) {\textbf{Cadre en Chêne}};
\end{tikzpicture}
\end{center}

\begin{enumerate}
    \item \capSApproprier\ Exprimer le périmètre $P(x)$ du cadre en fonction de la variable $x$ sous forme d'une somme :
    \[ P(x) = x + (x + 3) + x + (x + 3) \]
    \item \capRealiser\ Réduire cette expression en regroupant les termes en $x$ et les nombres :
    \lignesreponse[2]
    \item \capRealiser\ Exprimer l'aire $A(x)$ du rectangle en fonction de $x$ sous la forme d'un produit :
    \[ A(x) = x \times (x + 3) \]
    \item \capAnalyser\ Développer ce produit pour obtenir l'expression de l'aire sous forme d'une somme :
    \lignesreponse[2]
    \item \capValider\ Calculer le périmètre et l'aire pour une largeur $x = 2\text{ m}$ :
    \lignesreponse[2]
\end{enumerate}
\end{activite}

% ==============================================================================
% 3. COURS ET NOTIONS ESSENTIELLES
% ==============================================================================
\section{Cours : Conventions, Développement et Réduction}

\subsection{Conventions d'écriture et Réduction}

\begin{cadredef}[Convention de simplification]
Dans une expression algébrique, le signe de multiplication $\times$ peut être omis devant une lettre ou une parenthèse :
\[ 3 \times x = 3x \qquad 1 \times x = x \qquad x \times x = x^2 \qquad 5 \times (x + 2) = 5(x + 2) \]
\end{cadredef}

\begin{cadremethode}[Réduire une somme littérale]
Pour réduire une expression littérale, on regroupe et on additionne les termes de \textbf{même nature} (les termes en $x^2$ ensemble, les termes en $x$ ensemble, les constantes numériques ensemble).
\end{cadremethode}

\begin{cadreexemple}[Réduction guidée]
Réduire l'expression $A = 5x^2 + 3x - 7 - 2x^2 + 8x + 4$ :
\begin{align*}
    A &= (5x^2 - 2x^2) + (3x + 8x) + (-7 + 4) \\
    A &= 3x^2 + 11x - 3
\end{align*}
\end{cadreexemple}

\begin{cadrepiege}[Erreur classique d'addition]
On ne peut \textbf{jamais} additionner des termes de familles différentes !
\[ 3x + 4 \neq 7x \qquad\text{et}\qquad 2x^2 + 3x \neq 5x^3 \]
\end{cadrepiege}

\subsection{Développement : Distributivité Simple et Double}

\begin{cadreprop}[Formules de distributivité]
Pour tous nombres réels $k$, $a$, $b$, $c$ et $d$ :
\begin{itemize}
    \item \textbf{Distributivité Simple} :
    \[ k(a + b) = k \times a + k \times b = ka + kb \]
    \[ k(a - b) = k \times a - k \times b = ka - kb \]
    \item \textbf{Double Distributivité} :
    \[ (a + b)(c + d) = ac + ad + bc + bd \]
\end{itemize}
\end{cadreprop}

\begin{cadreexemple}[Développement d'une double distributivité]
Développer et réduire l'expression $B = (2x + 3)(x - 5)$ :
\begin{align*}
    B &= 2x \times x + 2x \times (-5) + 3 \times x + 3 \times (-5) \\
    B &= 2x^2 - 10x + 3x - 15 \\
    B &= 2x^2 - 7x - 15
\end{align*}
\end{cadreexemple}

\subsection{Factorisation et Identité Remarquable}

\begin{cadredef}[Factoriser]
\textbf{Factoriser} une expression consiste à transformer une somme (ou différence) en un \textbf{produit de facteurs}.
\[ ka + kb = k(a + b) \]
\end{cadredef}

\begin{cadreprop}[Identité remarquable : Différence de deux carrés]
Pour tous nombres réels $a$ et $b$ :
\[ a^2 - b^2 = (a - b)(a + b) \]
\end{cadreprop}

\begin{cadreexemple}[Factorisations usuelles]
\begin{itemize}
    \item Par facteur commun : $C = 6x + 18 = 6 \times x + 6 \times 3 = 6(x + 3)$
    \item Par identité remarquable : $D = x^2 - 49 = x^2 - 7^2 = (x - 7)(x + 7)$
\end{itemize}
\end{cadreexemple}

% ==============================================================================
% 4. AUTOMATISMES & QUESTIONS FLASH
% ==============================================================================
\section{Questions Flash / Automatismes}

\begin{automatisme}[Calcul Mental \& Réflexes Algébriques (10 min)]
\noindent
\begin{tabularx}{\linewidth}{X|X}
\textbf{1.} Réduire : $4x + 7x = \pointille[2cm]$ & \textbf{6.} Développer : $3(2x + 5) = \pointille[2.5cm]$ \\
\textbf{2.} Réduire : $5x^2 - 9x^2 = \pointille[2cm]$ & \textbf{7.} Développer : $-2(x - 4) = \pointille[2.5cm]$ \\
\textbf{3.} Calculer $A = 3x - 2$ pour $x = 4$ : & \textbf{8.} Factoriser : $7x + 21 = \pointille[2.5cm]$ \\
$\implies A = \pointille[2cm]$ & \textbf{9.} Factoriser : $x^2 - 16 = \pointille[2.5cm]$ \\
\textbf{4.} Calculer $B = x^2 + 5$ pour $x = -3$ : & \textbf{10.} Vrai ou Faux ? \\
$\implies B = \pointille[2cm]$ & $(x+2)^2 = x^2 + 4$ : \pointille[2cm] \\
\textbf{5.} Simplifier : $3x \times 4x = \pointille[2cm]$ & \textbf{Bonus :} Développer $(x+1)(x+3) = \pointille[2.5cm]$ \\
\end{tabularx}
\end{automatisme}

% ==============================================================================
% 5. TRAVAUX DIRIGÉS (TD) - 10 EXERCICES CORRIGÉS
% ==============================================================================
\section{Travaux Dirigés (TD) : Exercices d'Entraînement}

\begin{exercice}[2 pts]{\capRealiser}{Simplification et Réduction}
Réduire les expressions suivantes au maximum :
\begin{enumerate}
    \item $A = 7x + 3 - 2x + 9$
    \item $B = 4x^2 - 6x + 8 - x^2 + 10x - 12$
\end{enumerate}
\lignesreponse[4]
\end{exercice}

\begin{exercice}[2 pts]{\capRealiser}{Développement par Distributivité Simple}
Développer puis réduire les expressions :
\begin{enumerate}
    \item $C = 4(3x - 5)$
    \item $D = -3x(2x - 7)$
    \item $E = 5(x + 2) - 2(3x - 4)$
\end{enumerate}
\lignesreponse[6]
\end{exercice}

\begin{exercice}[3 pts]{\capRealiser}{Double Distributivité}
Développer, réduire et ordonner suivant les puissances décroissantes de $x$ :
\begin{enumerate}
    \item $F = (x + 4)(x + 7)$
    \item $G = (2x - 3)(x + 5)$
    \item $H = (3x - 2)(4x - 1)$
\end{enumerate}
\lignesreponse[6]
\end{exercice}

\begin{exercice}[2 pts]{\capRealiser}{Factorisation et Identités}
Factoriser chacune des expressions suivantes :
\begin{enumerate}
    \item $I = 5x + 35$
    \item $J = x^2 - 81$
    \item $K = 4x^2 - 25$
\end{enumerate}
\lignesreponse[5]
\end{exercice}

\begin{exercice}[3 pts]{\capAnalyser}{Problème Métier BTP : Périmètre d'une clôture}
Sur un chantier de construction, un chef d'équipe doit sécuriser une zone rectangulaire de longueur $L = 3x + 10$ mètres et de largeur $\ell = 2x + 4$ mètres.
\begin{enumerate}
    \item Exprimer le périmètre $P$ de la zone en fonction de $x$ sous forme développée et réduite.
    \item Calculer la longueur totale de grillage nécessaire si $x = 6\text{ m}$.
    \item Le grillage est vendu en rouleaux de $25\text{ m}$ à $48\text{ \euro}$ l'unité. Déterminer le coût total d'achat.
\end{enumerate}
\lignesreponse[6]
\end{exercice}

\begin{exercice}[3 pts]{\capAnalyser}{Problème Aménagement : Aire d'une terrasse}
Un paysagiste aménage une terrasse dont l'aire totale est donnée par la formule $A(x) = (x + 5)(2x + 3) - (x^2 + 15)$.
\begin{enumerate}
    \item Développer et réduire l'expression $A(x)$.
    \item Calculer l'aire de la terrasse pour $x = 4\text{ m}$.
\end{enumerate}
\lignesreponse[5]
\end{exercice}

% ==============================================================================
% 6. TP TICE TABLEUR
% ==============================================================================
\section{TP TICE : Modélisation et Tableur}

\begin{tpinfo}[Tableur LibreOffice Calc / Excel]{Formules et Équivalence Algébrique}
Un artisan teste deux tarifs d'installation électrique :
\begin{itemize}
    \item \textbf{Formule A} : $50\text{ \euro}$ de forfait déplacement $+ 35\text{ \euro}$ par heure : $T_A(x) = 35x + 50$
    \item \textbf{Formule B} : $20\text{ \euro}$ de forfait $+ 42\text{ \euro}$ par heure : $T_B(x) = 42x + 20$
\end{itemize}

\begin{center}
\begin{tabular}{|c|c|c|c|}
\hline
\rowcolor{slate200} & \textbf{A} & \textbf{B} & \textbf{C} \\
\hline
\textbf{1} & \textbf{Nombre d'heures ($x$)} & \textbf{Tarif A (\euro)} & \textbf{Tarif B (\euro)} \\
\hline
\textbf{2} & 1 & \texttt{=35*A2+50} & \texttt{=42*A2+20} \\
\hline
\textbf{3} & 2 & & \\
\hline
\textbf{4} & 3 & & \\
\hline
\textbf{5} & 4 & & \\
\hline
\textbf{6} & 5 & & \\
\hline
\end{tabular}
\end{center}

\begin{enumerate}
    \item Quelle formule tableur doit-on saisir dans la cellule \texttt{B2} puis étirer vers le bas ?
    \item À partir de combien d'heures de travail la Formule A devient-elle plus avantageuse pour le client ?
\end{enumerate}
\lignesreponse[3]
\end{tpinfo}

% ==============================================================================
% 7. VERS LE BREVET (DNB PRO)
% ==============================================================================
\section{Vers le Brevet (DNB Pro)}

\begin{sujetdnb}[DNB Pro Métropole][6 pts]{Menuiserie \& Terrasse en Bois}
Un artisan menuisier réalise une terrasse en lames de bois composite. Le plan ci-dessous indique les dimensions exprimées en mètres :

\begin{center}
\begin{tikzpicture}[scale=0.85]
    \draw[fill=colDNBBg, draw=colDNB, line width=1pt] (0,0) -- (7,0) -- (7,3) -- (4,3) -- (4,5) -- (0,5) -- cycle;
    \node at (2,2.5) {\textbf{Partie 1}};
    \node at (5.5,1.5) {\textbf{Partie 2}};
    \draw[<->, >=stealth, slate700] (0,-0.3) -- (7,-0.3) node[midway, below] {$x + 4$};
    \draw[<->, >=stealth, slate700] (-0.3,0) -- (-0.3,5) node[midway, left] {$x + 2$};
    \draw[<->, >=stealth, slate700] (7.3,0) -- (7.3,3) node[midway, right] {$3$};
    \draw[<->, >=stealth, slate700] (4,5.3) -- (0,5.3) node[midway, above] {$4$};
\end{tikzpicture}
\end{center}

\begin{enumerate}
    \item \capSApproprier\ Montrer que l'aire totale de la terrasse s'exprime par la relation :
    \[ \mathcal{A}(x) = 4(x + 2) + 3x \]
    \item \capRealiser\ Développer et réduire cette expression.
    \item \capValider\ Pour une terrasse réelle avec $x = 3\text{ m}$ :
    \begin{enumerate}
        \item Calculer l'aire exacte en $\text{m}^2$.
        \item Les lames coûtent $32\text{ \euro}/\text{m}^2$. Calculer le montant total du bois nécessaire.
    \end{enumerate}
\end{enumerate}
\lignesreponse[6]
\end{sujetdnb}

% ==============================================================================
% 8. CORRIGÉS DÉTAILLÉS
% ==============================================================================
\section{Corrigés Détaillés du Chapitre}

\begin{solution}[Corrigé des Exercices de TD]
\begin{itemize}
    \item \textbf{Exercice 1} :
    \begin{enumerate}
        \item $A = (7x - 2x) + (3 + 9) = \repbox{5x + 12}$
        \item $B = (4x^2 - x^2) + (-6x + 10x) + (8 - 12) = \repbox{3x^2 + 4x - 4}$
    \end{enumerate}
    \item \textbf{Exercice 2} :
    \begin{enumerate}
        \item $C = 4 \times 3x + 4 \times (-5) = \repbox{12x - 20}$
        \item $D = -3x \times 2x - 3x \times (-7) = \repbox{-6x^2 + 21x}$
        \item $E = 5x + 10 - 6x + 8 = \repbox{-x + 18}$
    \end{enumerate}
    \item \textbf{Exercice 3} :
    \begin{enumerate}
        \item $F = x^2 + 7x + 4x + 28 = \repbox{x^2 + 11x + 28}$
        \item $G = 2x^2 + 10x - 3x - 15 = \repbox{2x^2 + 7x - 15}$
        \item $H = 12x^2 - 3x - 8x + 2 = \repbox{12x^2 - 11x + 2}$
    \end{enumerate}
    \item \textbf{Exercice 4} :
    \begin{enumerate}
        \item $I = 5(x + 7)$
        \item $J = x^2 - 9^2 = (x - 9)(x + 9)$
        \item $K = (2x)^2 - 5^2 = (2x - 5)(2x + 5)$
    \end{enumerate}
    \item \textbf{Exercice 5 (BTP)} :
    \begin{enumerate}
        \item $P = 2(L + \ell) = 2((3x + 10) + (2x + 4)) = 2(5x + 14) = \repbox{10x + 28}\text{ m}$.
        \item Pour $x = 6$ : $P = 10 \times 6 + 28 = 60 + 28 = \repbox{88\text{ m}}$.
        \item Nombre de rouleaux : $88 / 25 = 3{,}52 \implies$ il faut \textbf{4 rouleaux}. Coût : $4 \times 48 = \repbox{192\text{ \euro}}$.
    \end{enumerate}
\end{itemize}
\end{solution}

\begin{solution}[Corrigé du Sujet DNB Pro]
\begin{enumerate}
    \item L'aire est la somme de l'aire du rectangle 1 ($4 \times (x + 2)$) et du rectangle 2 ($3 \times x$). Donc $\mathcal{A}(x) = 4(x+2) + 3x$.
    \item $\mathcal{A}(x) = 4x + 8 + 3x = \repbox{7x + 8}$.
    \item Pour $x = 3\text{ m}$ :
    \begin{enumerate}
        \item $\mathcal{A}(3) = 7 \times 3 + 8 = 21 + 8 = \repbox{29\text{ m}^2}$.
        \item Coût du bois : $29 \times 32 = \repbox{928\text{ \euro}}$.
    \end{enumerate}
\end{enumerate}
\end{solution}


\chapter{Équations du Premier Degré \& Problèmes}
\label{chap:equations}

% ==============================================================================
% 1. OBJECTIFS & COMPÉTENCES DU RÉFÉRENTIEL
% ==============================================================================
\begin{cadreretenu}[Objectifs \& Compétences du DNB Pro]
\begin{itemize}
    \item \textbf{Connaissances} : Notion d'égalité, d'inconnue, de membre d'une équation, équation produit nul.
    \item \textbf{Capacités} :
    \begin{itemize}
        \item \capRealiser\ Résoudre algébriquement une équation de la forme $ax + b = c$ et $ax + b = cx + d$.
        \item \capRealiser\ Résoudre une équation produit nul de type $(ax + b)(cx + d) = 0$.
        \item \capAnalyser\ Mettre en équation un problème issu de la vie courante ou d'un domaine professionnel (choix de l'inconnue, mise en équation, résolution, conclusion).
        \item \capValider\ Tester et vérifier si un nombre est solution d'une équation.
    \end{itemize}
\end{itemize}
\end{cadreretenu}

\vspace{0.3cm}

% ==============================================================================
% 2. ACTIVITÉ D'INVESTIGATION
% ==============================================================================
\section{Activité d'Introduction : Le Principe de la Balance}

\begin{activite}{L'équilibre d'une balance à deux plateaux}
Une balance est en parfait équilibre. Sur le plateau de gauche, on a placé 3 masses identiques inconnues notées $x$ et une masse marquée de $5\text{ kg}$. Sur le plateau de droite, on a placé une masse $x$ et une masse marquée de $17\text{ kg}$.

\begin{center}
\begin{tikzpicture}[scale=0.85]
    % Socle et fléau
    \draw[fill=slate700] (-0.3,0) rectangle (0.3,2);
    \draw[fill=slate800] (-1.5,0) rectangle (1.5,0.3);
    \draw[line width=2pt, colPrimary] (-4,2) -- (4,2);
    \draw[fill=colSecondary] (0,2) circle (4pt);

    % Plateaux
    \draw[thick, slate600] (-4,2) -- (-4,1) -- (-2.5,1) -- (-4,2);
    \draw[thick, slate600] (4,2) -- (4,1) -- (2.5,1) -- (4,2);
    \draw[line width=1.5pt, slate800] (-4.5,1) -- (-2.2,1);
    \draw[line width=1.5pt, slate800] (2.2,1) -- (4.5,1);

    % Masses gauche
    \draw[fill=colExoBg, draw=colExo, thick] (-4.2,1) rectangle (-3.6,1.6) node[midway, font=\footnotesize\bfseries] {$x$};
    \draw[fill=colExoBg, draw=colExo, thick] (-3.5,1) rectangle (-2.9,1.6) node[midway, font=\footnotesize\bfseries] {$x$};
    \draw[fill=colExoBg, draw=colExo, thick] (-3.85,1.6) rectangle (-3.25,2.2) node[midway, font=\footnotesize\bfseries] {$x$};
    \draw[fill=amber!30, draw=amber!80!black, thick] (-2.8,1) rectangle (-2.3,1.5) node[midway, font=\tiny\bfseries] {$5\text{kg}$};

    % Masses droite
    \draw[fill=colExoBg, draw=colExo, thick] (2.5,1) rectangle (3.1,1.6) node[midway, font=\footnotesize\bfseries] {$x$};
    \draw[fill=amber!30, draw=amber!80!black, thick] (3.3,1) rectangle (4.2,1.9) node[midway, font=\footnotesize\bfseries] {$17\text{kg}$};
\end{tikzpicture}
\end{center}

\begin{enumerate}
    \item \capSApproprier\ Traduire l'équilibre de cette balance par une équation mathématique :
    \[ 3x + 5 = x + 17 \]
    \item \capAnalyser\ Si on retire une masse $x$ sur chaque plateau, que devient l'égalité ?
    \lignesreponse[2]
    \item \capRealiser\ Si on retire ensuite $5\text{ kg}$ sur chaque plateau, quelle masse pèsent les deux masses $x$ restantes ? En déduire la valeur de la masse $x$ :
    \lignesreponse[3]
\end{enumerate}
\end{activite}

% ==============================================================================
% 3. COURS ET NOTIONS ESSENTIELLES
% ==============================================================================
\section{Cours : Résolution Algébrique et Méthodes}

\subsection{Définitions et Propriétés Fondamentales}

\begin{cadredef}[Équation du premier degré à une inconnue]
Une \textbf{équation du premier degré} est une égalité comportant une lettre inconnue (souvent $x$) élevée à la puissance 1. \\
\textbf{Résoudre} une équation, c'est trouver toutes les valeurs de $x$ qui rendent l'égalité vraie. Chaque valeur trouvée est appelée une \textbf{solution}.
\end{cadredef}

\begin{cadreprop}[Règles de conservation de l'égalité]
Une égalité reste vraie si :
\begin{itemize}
    \item On \textbf{ajoute} ou on \textbf{soustrait} un même nombre aux deux membres :
    \[ A = B \ssi A + c = B + c \quad\text{et}\quad A - c = B - c \]
    \item On \textbf{multiplie} ou on \textbf{divise} les deux membres par un même nombre \textbf{non nul} ($c \neq 0$) :
    \[ A = B \ssi A \times c = B \times c \quad\text{et}\quad \frac{A}{c} = \frac{B}{c} \]
\end{itemize}
\end{cadreprop}

\subsection{Méthode Générale de Résolution}

\begin{cadremethode}[Résoudre une équation du type $ax + b = cx + d$]
\begin{enumerate}
    \item \textbf{Regrouper} tous les termes contenant $x$ dans le membre de gauche (en soustrayant $cx$).
    \item \textbf{Regrouper} toutes les constantes numériques dans le membre de droite (en soustrayant $b$).
    \item \textbf{Réduire} chaque membre pour obtenir une équation simplifiée de la forme $k x = m$.
    \item \textbf{Diviser} les deux membres par le coefficient $k$ ($k \neq 0$) : $x = \frac{m}{k}$.
    \item \textbf{Conclure} en vérifiant la solution trouvée.
\end{enumerate}
\end{cadremethode}

\begin{cadreexemple}[Résolution pas à pas]
Résoudre dans $\mathbb{R}$ l'équation $5x - 8 = 2x + 7$ :
\begin{align*}
    5x - 8 - 2x &= 7 && \text{(on soustrait $2x$ des deux côtés)} \\
    3x - 8 &= 7 && \text{(on réduit le membre de gauche)} \\
    3x &= 7 + 8 && \text{(on ajoute $8$ des deux côtés)} \\
    3x &= 15 \\
    x &= \frac{15}{3} = \repbox{5} && \text{(on divise par $3$)}
\end{align*}
\textbf{Vérification} : Pour $x = 5$, Membre de gauche $= 5 \times 5 - 8 = 17$, Membre de droite $= 2 \times 5 + 7 = 17$. L'égalité est vérifiée, la solution unique est bien $x = 5$.
\end{cadreexemple}

\subsection{Équations Produits Nuls}

\begin{cadreprop}[Règle du produit nul]
Un produit de facteurs est nul si et seulement si \textbf{au moins l'un des facteurs est nul} :
\[ A \times B = 0 \ssi A = 0 \quad\text{ou}\quad B = 0 \]
\end{cadreprop}

\begin{cadreexemple}[Résolution d'une équation produit]
Résoudre $(2x - 6)(3x + 12) = 0$ :
\begin{align*}
    2x - 6 = 0 \quad &\text{ou} \quad 3x + 12 = 0 \\
    2x = 6 \quad &\text{ou} \quad 3x = -12 \\
    x = \frac{6}{2} = \repbox{3} \quad &\text{ou} \quad x = \frac{-12}{3} = \repbox{-4}
\end{align*}
L'équation admet donc deux solutions : $x_1 = 3$ et $x_2 = -4$.
\end{cadreexemple}

% ==============================================================================
% 4. AUTOMATISMES & QUESTIONS FLASH
% ==============================================================================
\section{Questions Flash / Automatismes}

\begin{automatisme}[Résolutions Mentales \& Réflexes (10 min)]
\noindent
\begin{tabularx}{\linewidth}{X|X}
\textbf{1.} $x + 7 = 15 \implies x = \pointille[2cm]$ & \textbf{6.} $3x + 4 = 19 \implies x = \pointille[2cm]$ \\
\textbf{2.} $x - 9 = 4 \implies x = \pointille[2cm]$ & \textbf{7.} $5x - 10 = 0 \implies x = \pointille[2cm]$ \\
\textbf{3.} $4x = 28 \implies x = \pointille[2cm]$ & \textbf{8.} $2x = x + 8 \implies x = \pointille[2cm]$ \\
\textbf{4.} $-3x = 18 \implies x = \pointille[2cm]$ & \textbf{9.} $(x - 5)(x + 2) = 0 \implies x = \pointille[1.5cm] \text{ ou } \pointille[1.5cm]$ \\
\textbf{5.} $\frac{x}{5} = 6 \implies x = \pointille[2cm]$ & \textbf{10.} $x = 3$ est-il solution de $2x + 1 = 7$ ? \pointille[1.5cm] \\
\end{tabularx}
\end{automatisme}

% ==============================================================================
% 5. TRAVAUX DIRIGÉS (TD) - 10 EXERCICES CORRIGÉS
% ==============================================================================
\section{Travaux Dirigés (TD) : Exercices d'Entraînement}

\begin{exercice}[2 pts]{\capRealiser}{Équations de base à une étape}
Résoudre les équations suivantes :
\begin{enumerate}
    \item $x + 14 = 31$
    \item $y - 8{,}5 = 12$
    \item $6x = 42$
    \item $-7t = 56$
\end{enumerate}
\lignesreponse[4]
\end{exercice}

\begin{exercice}[3 pts]{\capRealiser}{Équations du type $ax + b = c$}
Résoudre dans $\mathbb{R}$ :
\begin{enumerate}
    \item $3x + 7 = 25$
    \item $5x - 13 = 17$
    \item $-4x + 9 = -11$
    \item $2x + 15 = 4$
\end{enumerate}
\lignesreponse[6]
\end{exercice}

\begin{exercice}[3 pts]{\capRealiser}{Équations avec l'inconnue dans les deux membres}
Résoudre les équations :
\begin{enumerate}
    \item $7x - 4 = 3x + 16$
    \item $8x + 5 = 2x - 19$
    \item $2(3x - 1) = 4x + 10$
\end{enumerate}
\lignesreponse[6]
\end{exercice}

\begin{exercice}[2 pts]{\capRealiser}{Équations Produits Nuls}
Résoudre les équations produit :
\begin{enumerate}
    \item $(x - 7)(2x + 8) = 0$
    \item $(3x - 15)(4x + 2) = 0$
\end{enumerate}
\lignesreponse[4]
\end{exercice}

\begin{exercice}[3 pts]{\capAnalyser}{Problème Plomberie : Facturation de dépannage}
Un plombier applique un forfait déplacement de $45\text{ \euro}$ et un tarif horaire de $38\text{ \euro}$ de l'heure.
Une facture totale s'élève à $159\text{ \euro}$ TTC.
\begin{enumerate}
    \item Écrire une équation traduisant cette situation en notant $h$ le nombre d'heures de main d'œuvre.
    \item Résoudre cette équation et indiquer la durée exacte de l'intervention.
\end{enumerate}
\lignesreponse[4]
\end{exercice}

\begin{exercice}[3 pts]{\capAnalyser}{Problème Mécanique : Équilibre de charges}
Un atelier automobile charge une camionnette. La masse à vide du véhicule est de $1\,450\text{ kg}$. On y charge 12 caisses d'outillage de même masse $m$ et un groupe électrogène de $130\text{ kg}$.
La masse totale du véhicule chargé est mesurée à $2\,300\text{ kg}$.
\begin{enumerate}
    \item Poser l'équation permettant de calculer la masse $m$ d'une caisse.
    \item Résoudre cette équation et conclure.
\end{enumerate}
\lignesreponse[5]
\end{exercice}

% ==============================================================================
% 6. SUJET TYPE DNB PRO
% ==============================================================================
\section{Vers le Brevet (DNB Pro)}

\begin{sujetdnb}[DNB Pro Centres Étrangers][8 pts]{Location de Véhicules Utilitaires}
Pour transporter des matériaux de chantier, une entreprise hésite entre deux agences de location d'utilitaires pour une journée :
\begin{itemize}
    \item \textbf{Agence RapideLoc} : Forfait journalier de $60\text{ \euro}$ plus $0{,}25\text{ \euro}$ par kilomètre parcouru.
    \item \textbf{Agence EcoDrive} : Forfait journalier de $30\text{ \euro}$ plus $0{,}40\text{ \euro}$ par kilomètre parcouru.
\end{itemize}
On note $x$ la distance parcourue en kilomètres au cours de la journée.

\begin{enumerate}
    \item \capSApproprier\ Exprimer en fonction de $x$ le coût de location $C_1(x)$ pour RapideLoc et $C_2(x)$ pour EcoDrive.
    \item \capRealiser\ Calculer le coût pour chaque agence si l'entreprise parcourt une distance de $100\text{ km}$.
    \item \capAnalyser\ Résoudre l'équation $60 + 0{,}25x = 30 + 0{,}40x$.
    \item \capValider\ Interpréter la valeur de $x$ trouvée dans le contexte commercial : pour quelle distance les deux agences facturent-elles exactement le même montant ?
    \item \capCommuniquer\ Conseiller l'entreprise si elle prévoit d'effectuer un trajet de $350\text{ km}$.
\end{enumerate}
\lignesreponse[7]
\end{sujetdnb}

% ==============================================================================
% 7. CORRIGÉS DÉTAILLÉS
% ==============================================================================
\section{Corrigés Détaillés du Chapitre}

\begin{solution}[Corrigé des Exercices de TD]
\begin{itemize}
    \item \textbf{Exercice 1} :
    1. $x = 31 - 14 = \repbox{17}$ \quad 2. $y = 12 + 8{,}5 = \repbox{20{,}5}$ \quad 3. $x = 42/6 = \repbox{7}$ \quad 4. $t = 56/(-7) = \repbox{-8}$.
    \item \textbf{Exercice 2} :
    1. $3x = 25 - 7 = 18 \implies x = \repbox{6}$ \quad 2. $5x = 30 \implies x = \repbox{6}$ \\
    3. $-4x = -20 \implies x = \repbox{5}$ \quad 4. $2x = 4 - 15 = -11 \implies x = \repbox{-5{,}5}$.
    \item \textbf{Exercice 3} :
    1. $7x - 3x = 16 + 4 \implies 4x = 20 \implies x = \repbox{5}$ \\
    2. $8x - 2x = -19 - 5 \implies 6x = -24 \implies x = \repbox{-4}$ \\
    3. $6x - 2 = 4x + 10 \implies 2x = 12 \implies x = \repbox{6}$.
    \item \textbf{Exercice 4} :
    1. $x = 7$ ou $2x = -8 \implies x = \repbox{-4}$. Solutions : $\{ -4 ; 7 \}$. \\
    2. $3x = 15 \implies x = \repbox{5}$ ou $4x = -2 \implies x = \repbox{-0{,}5}$. Solutions : $\{ -0{,}5 ; 5 \}$.
    \item \textbf{Exercice 5 (Plomberie)} :
    1. Équation : $38h + 45 = 159$. \\
    2. $38h = 159 - 45 = 114 \implies h = 114 / 38 = \repbox{3\text{ heures}}$.
    \item \textbf{Exercice 6 (Mécanique)} :
    1. Équation : $1450 + 12m + 130 = 2300 \implies 12m + 1580 = 2300$. \\
    2. $12m = 2300 - 1580 = 720 \implies m = 720 / 12 = \repbox{60\text{ kg}}$.
\end{itemize}
\end{solution}

\begin{solution}[Corrigé du Sujet DNB Pro]
\begin{enumerate}
    \item $C_1(x) = 0{,}25x + 60$ et $C_2(x) = 0{,}40x + 30$.
    \item Pour $x = 100\text{ km}$ :
    $C_1(100) = 0{,}25 \times 100 + 60 = 25 + 60 = 85\text{ \euro}$. \\
    $C_2(100) = 0{,}40 \times 100 + 30 = 40 + 30 = 70\text{ \euro}$. EcoDrive est la moins chère.
    \item $60 - 30 = 0{,}40x - 0{,}25x \implies 30 = 0{,}15x \implies x = \frac{30}{0{,}15} = \repbox{200}$.
    \item Pour une distance de \textbf{200 km}, les deux agences facturent exactement le même tarif ($110\text{ \euro}$).
    \item Pour $x = 350\text{ km} > 200\text{ km}$, c'est l'agence **RapideLoc** qui est la plus économique ($C_1(350) = 147{,}50\text{ \euro}$ contre $C_2(350) = 170\text{ \euro}$).
\end{enumerate}
\end{solution}


\chapter{Fonctions Linéaires, Affines \& Graphiques}
\label{chap:fonctions_affines}

% ==============================================================================
% 1. OBJECTIFS & COMPÉTENCES DU RÉFÉRENTIEL
% ==============================================================================
\begin{cadreretenu}[Objectifs \& Compétences du DNB Pro]
\begin{itemize}
    \item \textbf{Connaissances} : Notion de fonction, variable, image, antécédent, fonction linéaire $f(x)=ax$, fonction affine $f(x)=ax+b$, coefficient directeur, ordonnée à l'origine.
    \item \textbf{Capacités} :
    \begin{itemize}
        \item \capRealiser\ Calculer l'image d'un nombre par une fonction définie par une formule.
        \item \capRealiser\ Déterminer un antécédent en résolvant une équation du premier degré.
        \item \capRealiser\ Représenter graphiquement une fonction linéaire (droite passant par l'origine) ou affine.
        \item \capAnalyser\ Lire graphiquement des images et des antécédents, déterminer le coefficient directeur et l'ordonnée à l'origine.
        \item \capValider\ Modéliser une situation de proportionnalité ou d'abonnement / coût fixe par une fonction.
    \end{itemize}
\end{itemize}
\end{cadreretenu}

\vspace{0.3cm}

% ==============================================================================
% 2. COURS ET NOTIONS ESSENTIELLES
% ==============================================================================
\section{Cours : Notions Fondamentales \& Représentations}

\subsection{Vocabulaire des Fonctions : Image et Antécédent}

\begin{cadredef}[Fonction, Image, Antécédent]
Une \textbf{fonction} $f$ est un processus mathématique qui, à un nombre $x$ de départ (l'antécédent), associe un unique nombre d'arrivée noté $f(x)$ (l'image).
\[ x \overset{f}{\longmapsto} f(x) \]
\begin{itemize}
    \item Un nombre possède \textbf{une seule image}.
    \item Un nombre peut posséder \textbf{aucun, un ou plusieurs antécédents}.
\end{itemize}
\end{cadredef}

\begin{cadreexemple}[Calcul d'image et d'antécédent]
Soit la fonction $f$ définie par $f(x) = 3x - 5$.
\begin{itemize}
    \item \textbf{Calculer l'image de 4} : On remplace $x$ par $4$ : $f(4) = 3 \times 4 - 5 = 12 - 5 = \repbox{7}$.
    \item \textbf{Déterminer l'antécédent de 16} : On résout $f(x) = 16 \iff 3x - 5 = 16 \iff 3x = 21 \iff x = \repbox{7}$.
\end{itemize}
\end{cadreexemple}

\subsection{Fonctions Linéaires \& Fonctions Affines}

\begin{cadreprop}[Caractérisation des fonctions usuelles]
\begin{itemize}
    \item \textbf{Fonction Linéaire} : $f(x) = ax$. Sa courbe représentative est une \textbf{droite qui passe par l'origine} du repère $(0;0)$. Elle modélise une situation de \textbf{proportionnalité}.
    \item \textbf{Fonction Affine} : $f(x) = ax + b$. Sa courbe représentative est une \textbf{droite} qui coupe l'axe des ordonnées au point $(0\,;\,b)$.
    \begin{itemize}
        \item $a$ est le \textbf{coefficient directeur} (ou pente de la droite). Si $a > 0$, la fonction est croissante ; si $a < 0$, elle est décroissante.
        \item $b$ est l'\textbf{ordonnée à l'origine}.
    \end{itemize}
\end{itemize}
\end{cadreprop}

\begin{center}
\begin{tikzpicture}[scale=0.85]
    % Grille et repère
    \draw[step=1cm, slate200, very thin] (-3,-2) grid (5,4);
    \draw[thick, ->, >=stealth, slate700] (-3,0) -- (5,0) node[right] {$x$};
    \draw[thick, ->, >=stealth, slate700] (0,-2) -- (0,4) node[above] {$y$};
    \node[below left, font=\footnotesize] at (0,0) {$O$};

    % Droite affine f(x) = 0.8x + 1
    \draw[line width=1.4pt, colPrimary] (-3,-1.4) -- (3.75,4) node[above right] {$d_1 : f(x) = 0{,}8x + 1$};
    \fill[colPrimary] (0,1) circle (2.5pt) node[left=2pt, font=\footnotesize\bfseries] {$b = 1$};

    % Pente triangle
    \draw[thick, colSecondary, dashed] (0,1) -- (2,1) node[midway, below] {$+2$};
    \draw[thick, colSecondary, dashed] (2,1) -- (2,2.6) node[midway, right] {$+1{,}6$};

    % Droite linéaire g(x) = -x
    \draw[line width=1.4pt, colProp] (-3,3) -- (3,-3) node[below right] {$d_2 : g(x) = -x$};
    \fill[colProp] (0,0) circle (2.5pt);
\end{tikzpicture}
\end{center}

% ==============================================================================
% 3. AUTOMATISMES & QUESTIONS FLASH
% ==============================================================================
\section{Questions Flash / Automatismes}

\begin{automatisme}[Lecture Graphique \& Calculs Rapides (10 min)]
Soit la fonction affine $g(x) = -2x + 6$.
\noindent
\begin{tabularx}{\linewidth}{X|X}
\textbf{1.} Le coefficient directeur est $a = \pointille[1.5cm]$ & \textbf{4.} L'image de $3$ est $g(3) = \pointille[2cm]$ \\
\textbf{2.} L'ordonnée à l'origine est $b = \pointille[1.5cm]$ & \textbf{5.} L'antécédent de $0$ est : $x = \pointille[2cm]$ \\
\textbf{3.} La fonction $g$ est : (Croissante / Décroissante) ? \pointille[2cm] & \textbf{6.} La droite passe par le point $(0\,;\,\pointille[1cm])$. \\
\end{tabularx}
\end{automatisme}

% ==============================================================================
% 4. TRAVAUX DIRIGÉS (TD)
% ==============================================================================
\section{Travaux Dirigés (TD) : Exercices d'Entraînement}

\begin{exercice}[3 pts]{\capRealiser}{Calculs d'images et d'antécédents}
On considère la fonction $f$ définie sur $\mathbb{R}$ par $f(x) = 4x - 7$.
\begin{enumerate}
    \item Calculer $f(2)$, $f(-3)$ et $f(0)$.
    \item Déterminer l'antécédent de $13$ par la fonction $f$.
    \item Le point $M(5\,;\,13)$ appartient-il à la droite représentative de $f$ ? Justifier.
\end{enumerate}
\lignesreponse[5]
\end{exercice}

\begin{exercice}[4 pts]{\capRealiser}{Tracé de droites représentatives}
Dans un repère orthonormé d'unité $1\text{ cm}$ pour 1 unité sur chaque axe :
\begin{enumerate}
    \item Construire le tableau de valeurs pour la fonction $f(x) = 2x - 3$ avec $x \in \{-1 ; 0 ; 2 ; 4\}$.
    \item Tracer la droite $(d_1)$ représentative de $f$.
    \item Tracer sur le même repère la droite $(d_2)$ représentative de la fonction linéaire $g(x) = -1{,}5x$.
\end{enumerate}
\quadrillage[4.5cm]
\end{exercice}

\begin{exercice}[4 pts]{\capAnalyser}{Problème Commercial : Forfaits de Téléphonie}
Un opérateur propose deux forfaits pour une flotte d'entreprise :
\begin{itemize}
    \item \textbf{Offre A} : Sans abonnement, $0{,}15\euro$ la minute d'appel : $A(x) = 0{,}15x$.
    \item \textbf{Offre B} : Abonnement fixe de $12\euro$ par mois puis $0{,}05\euro$ la minute : $B(x) = 0{,}05x + 12$.
\end{itemize}
\begin{enumerate}
    \item Préciser la nature de chaque fonction (linéaire ou affine).
    \item Calculer le montant de la facture pour 60 minutes de communication avec chaque offre.
    \item Résoudre l'équation $0{,}15x = 0{,}05x + 12$ et interpréter le résultat.
\end{enumerate}
\lignesreponse[5]
\end{exercice}

% ==============================================================================
% 5. VERS LE BREVET (DNB PRO)
% ==============================================================================
\section{Vers le Brevet (DNB Pro)}

\begin{sujetdnb}[DNB Pro Métropole][10 pts]{Gestion d'un Club Nautique}
Un club de kayak propose deux tarifs de location de matériel pour la saison d'été :
\begin{itemize}
    \item \textbf{Tarif Plein} : $18\euro$ par séance de 2 heures.
    \item \textbf{Tarif Club (avec adhésion)} : Carte annuelle à $40\euro$ puis $10\euro$ par séance de 2 heures.
\end{itemize}
On note $x$ le nombre de séances effectuées par un adhérent.

\begin{center}
\begin{tikzpicture}[scale=0.75]
    \draw[step=1cm, slate200, very thin] (0,0) grid (10,10);
    \draw[thick, ->, >=stealth, slate700] (0,0) -- (10.5,0) node[right] {Nombre de séances $x$};
    \draw[thick, ->, >=stealth, slate700] (0,0) -- (0,10.5) node[above] {Prix total (\euro)};

    % Graduations (axe x : 1 car = 1 seance, axe y : 1 car = 20 euros)
    \foreach \x in {1,2,...,10} \draw (\x,0.1) -- (\x,-0.1) node[below, font=\scriptsize] {\x};
    \foreach \y in {20,40,...,200} \draw (0.1,\y/20) -- (-0.1,\y/20) node[left, font=\scriptsize] {\y};

    % Droites
    % Tarif Plein: 18x -> y = 18*x/20
    \draw[line width=1.3pt, colPrimary] (0,0) -- (10,9) node[above right] {$T_1(x) = 18x$};
    % Tarif Club: 10x + 40 -> y = (10*x + 40)/20 = 0.5x + 2
    \draw[line width=1.3pt, colExem] (0,2) -- (10,7) node[above right] {$T_2(x) = 10x + 40$};

    \fill[red] (5,4.5) circle (3pt) node[above left, text=red, font=\footnotesize\bfseries] {Point d'intersection};
\end{tikzpicture}
\end{center}

\begin{enumerate}
    \item \capSApproprier\ Associer à chaque tarif sa fonction représentative $T_1(x)$ et $T_2(x)$.
    \item \capAnalyser\ Par lecture graphique :
    \begin{enumerate}
        \item Quel est le prix à payer pour 3 séances au tarif plein ?
        \item Pour quel nombre de séances les deux tarifs sont-ils rigoureusement identiques ? Quel est alors le prix ?
        \item Pour un jeune qui prévoit de faire 8 séances dans l'été, quel est le tarif le plus avantageux ? Quelle sera son économie ?
    \end{enumerate}
    \item \capRealiser\ Retrouver par le calcul le nombre de séances pour lequel $T_1(x) = T_2(x)$.
\end{enumerate}
\lignesreponse[7]
\end{sujetdnb}

% ==============================================================================
% 6. CORRIGÉS DÉTAILLÉS
% ==============================================================================
\section{Corrigés Détaillés du Chapitre}

\begin{solution}[Corrigé des Exercices et du Sujet DNB]
\textbf{Exercice 1} :
1. $f(2) = 4(2)-7 = 1$ ; $f(-3) = 4(-3)-7 = -19$ ; $f(0) = -7$. \\
2. $4x - 7 = 13 \implies 4x = 20 \implies x = \repbox{5}$. \\
3. $f(5) = 4(5)-7 = 13$. Donc le point $M(5;13)$ \textbf{appartient bien} à la droite.

\vspace{2mm}
\textbf{Exercice 3 (Téléphonie)} :
1. $A$ est une fonction linéaire (proportionnalité), $B$ est une fonction affine. \\
2. $A(60) = 0{,}15 \times 60 = 9\euro$ ; $B(60) = 0{,}05 \times 60 + 12 = 15\euro$. \\
3. $0{,}15x - 0{,}05x = 12 \implies 0{,}10x = 12 \implies x = \repbox{120\text{ minutes}}$. Au-delà de 2 heures d'appel, l'Offre B est plus rentable.

\vspace{2mm}
\textbf{Sujet DNB Pro (Club Nautique)} :
1. $T_1(x) = 18x$ (Tarif Plein) et $T_2(x) = 10x + 40$ (Tarif Club). \\
2. (a) Pour 3 séances au tarif plein : $18 \times 3 = 54\euro$ (graphiquement ordonnée 2,7). \\
(b) L'intersection a lieu pour \textbf{5 séances}, le montant est de \textbf{90 \euro}. \\
(c) Pour 8 séances : $T_1(8) = 144\euro$ et $T_2(8) = 120\euro$. Le **Tarif Club** est plus avantageux, économie de $144 - 120 = \repbox{24\euro}$. \\
3. $18x = 10x + 40 \implies 8x = 40 \implies x = \frac{40}{8} = \repbox{5\text{ séances}}$.
\end{solution}


\chapter{Théorème de Pythagore \& Trigonométrie}
\label{chap:pythagore_trigo}

% ==============================================================================
% 1. OBJECTIFS & COMPÉTENCES DU RÉFÉRENTIEL
% ==============================================================================
\begin{cadreretenu}[Objectifs \& Compétences du DNB Pro]
\begin{itemize}
    \item \textbf{Connaissances} : Triangle rectangle, hypoténuse, théorème direct de Pythagore, réciproque et contraposée, formules de trigonométrie ($\cos, \sin, \tan$).
    \item \textbf{Capacités} :
    \begin{itemize}
        \item \capRealiser\ Calculer la longueur de l'hypoténuse ou d'un côté de l'angle droit avec le théorème de Pythagore.
        \item \capAnalyser\ Démontrer si un triangle est rectangle ou non en utilisant la réciproque ou la contraposée.
        \item \capRealiser\ Calculer la mesure d'un angle aigu à l'aide des touches $\arccos, \arcsin, \arctan$ de la calculatrice.
        \item \capRealiser\ Calculer une longueur inconnue à l'aide des rapports trigonométriques.
        \item \capValider\ Résoudre un problème géométrique de charpente, menuiserie, pente de toit ou accessibilité PMR.
    \end{itemize}
\end{itemize}
\end{cadreretenu}

\vspace{0.3cm}

% ==============================================================================
% 2. COURS ET NOTIONS ESSENTIELLES
% ==============================================================================
\section{Cours : Théorème de Pythagore}

\subsection{Énoncé Direct et Calcul de Longueurs}

\begin{cadreprop}[Théorème de Pythagore]
Dans un triangle $ABC$ \textbf{rectangle en $A$}, le carré de la longueur de l'hypoténuse est égal à la somme des carrés des longueurs des deux autres côtés :
\[ BC^2 = AB^2 + AC^2 \]
\end{cadreprop}

\begin{center}
\begin{tikzpicture}[scale=0.9]
    \coordinate (A) at (0,0);
    \coordinate (B) at (4,0);
    \coordinate (C) at (0,3);

    \draw[fill=colDefBg, draw=colDef, line width=1.2pt] (A) -- (B) node[midway, below] {Côté $= 4\text{ cm}$} -- (C) node[midway, above right] {Hypoténuse $BC = 5\text{ cm}$} -- cycle node[midway, left] {Côté $= 3\text{ cm}$};

    % Angle droit
    \draw[thick, slate800] (0,0.3) -- (0.3,0.3) -- (0.3,0);
    \node[below left] at (A) {$A$};
    \node[below right] at (B) {$B$};
    \node[above left] at (C) {$C$};
\end{tikzpicture}
\end{center}

\begin{cadremethode}[Règles de calcul de longueurs]
\begin{itemize}
    \item \textbf{Pour calculer l'hypoténuse $BC$} :
    $BC^2 = AB^2 + AC^2 \implies BC = \sqrt{AB^2 + AC^2}$.
    \item \textbf{Pour calculer un côté de l'angle droit $AB$} :
    $AB^2 = BC^2 - AC^2 \implies AB = \sqrt{BC^2 - AC^2}$.
\end{itemize}
\end{cadremethode}

\subsection{Réciproque et Contraposée : Démontrer qu'un triangle est rectangle}

\begin{cadreprop}[Réciproque de Pythagore]
Dans un triangle, si le carré du plus grand côté est égal à la somme des carrés des deux autres côtés ($BC^2 = AB^2 + AC^2$), alors ce triangle est \textbf{rectangle} en $A$.
\end{cadreprop}

\begin{cadrepiege}[Rédaction rigoureuse exigée au DNB]
Au Brevet, il est \textbf{strictement interdit} d'écrire l'égalité dès le début ! On doit séparer les calculs :
\begin{enumerate}
    \item On calcule d'une part le carré du plus grand côté : $BC^2 = \dots$
    \item On calcule d'autre part la somme des carrés des deux autres côtés : $AB^2 + AC^2 = \dots$
    \item On compare les deux résultats et on conclut avec la propriété adéquate.
\end{enumerate}
\end{cadrepiege}

% ==============================================================================
% 3. COURS : TRIGONOMÉTRIE DANS LE TRIANGLE RECTANGLE
% ==============================================================================
\section{Cours : Trigonométrie}

\begin{cadredef}[Cosinus, Sinus, Tangente]
Dans un triangle rectangle, pour un angle aigu $\widehat{B}$ :
\begin{itemize}
    \item $\cos(\widehat{B}) = \dfrac{\text{Côté Adjacent}}{\text{Hypoténuse}} = \dfrac{AB}{BC}$
    \item $\sin(\widehat{B}) = \dfrac{\text{Côté Opposé}}{\text{Hypoténuse}} = \dfrac{AC}{BC}$
    \item $\tan(\widehat{B}) = \dfrac{\text{Côté Opposé}}{\text{Côté Adjacent}} = \dfrac{AC}{AB}$
\end{itemize}
\end{cadredef}

\begin{cadreretenu}[Moyen Mnémotechnique Incontournable : \textbf{CAH-SOH-TOA}]
\begin{center}
\textbf{C}osinus = \textbf{A}djacent / \textbf{H}ypoténuse \quad $\vert$ \quad
\textbf{S}inus = \textbf{O}pposé / \textbf{H}ypoténuse \quad $\vert$ \quad
\textbf{T}angente = \textbf{O}pposé / \textbf{A}djacent
\end{center}
\end{cadreretenu}

% ==============================================================================
% 4. AUTOMATISMES & QUESTIONS FLASH
% ==============================================================================
\section{Questions Flash / Automatismes}

\begin{automatisme}[Calculs Rapides \& Règle des Carrés (10 min)]
\noindent
\begin{tabularx}{\linewidth}{X|X}
\textbf{1.} $3^2 + 4^2 = \pointille[1.5cm] \implies \sqrt{\dots} = \pointille[1cm]$ & \textbf{4.} Dans le triangle rectangle, le plus grand côté s'appelle : \pointille[2.5cm] \\
\textbf{2.} $10^2 - 8^2 = \pointille[1.5cm] \implies \sqrt{\dots} = \pointille[1cm]$ & \textbf{5.} $\cos(60\degre) = \pointille[1.5cm]$ (valeur exacte ou décimale) \\
\textbf{3.} Les triplets pythagoriciens usuels : $(3, 4, 5)$ et $(6, 8, \pointille[1cm])$ & \textbf{6.} Sur la calculatrice, pour trouver l'angle dont le cosinus vaut 0,5, on tape : \pointille[2cm] \\
\end{tabularx}
\end{automatisme}

% ==============================================================================
% 5. TRAVAUX DIRIGÉS (TD)
% ==============================================================================
\section{Travaux Dirigés (TD) : Exercices d'Entraînement}

\begin{exercice}[3 pts]{\capRealiser}{Calcul direct de longueur}
Soit $RST$ un triangle rectangle en $S$ tel que $RS = 4{,}8\text{ cm}$ et $ST = 6{,}4\text{ cm}$.
\begin{enumerate}
    \item Faire un schéma à main levée.
    \item Calculer la longueur exacte de l'hypoténuse $RT$.
\end{enumerate}
\lignesreponse[4]
\end{exercice}

\begin{exercice}[3 pts]{\capAnalyser}{Démontrer si un triangle est rectangle}
Un triangle $KLM$ a pour dimensions : $KL = 7\text{ cm}$, $LM = 24\text{ cm}$ et $KM = 25\text{ cm}$.
Le triangle $KLM$ est-il rectangle ? Rédiger avec rigueur.
\lignesreponse[4]
\end{exercice}

\begin{exercice}[4 pts]{\capRealiser}{Calcul d'un angle avec la trigonométrie}
Dans un triangle $ABC$ rectangle en $A$, on donne $AB = 5\text{ cm}$ et $BC = 8\text{ cm}$.
\begin{enumerate}
    \item Écrire la relation trigonométrique donnant $\cos(\widehat{ABC})$.
    \item En déduire la mesure de l'angle $\widehat{ABC}$ arrondie au degré près.
    \item Calculer alors la mesure de l'angle $\widehat{ACB}$.
\end{enumerate}
\lignesreponse[4]
\end{exercice}

% ==============================================================================
% 6. VERS LE BREVET (DNB PRO)
% ==============================================================================
\section{Vers le Brevet (DNB Pro)}

\begin{sujetdnb}[DNB Pro Polynésie][9 pts]{Charpente \& Pente de Toiture}
Un charpentier pose les fermettes d'un comble de maison. La figure ci-dessous représente la coupe d'une ferme de charpente :

\begin{center}
\begin{tikzpicture}[scale=0.9]
    % Triangle principal
    \coordinate (A) at (0,0);
    \coordinate (B) at (8,0);
    \coordinate (S) at (4,2.5);
    \coordinate (H) at (4,0);

    \draw[fill=amber!15, draw=amber!80!black, line width=1.5pt] (A) -- (B) -- (S) -- cycle;
    \draw[dashed, thick, slate700] (S) -- (H) node[midway, right] {Poinçon $SH = 2{,}4\text{ m}$};

    % Angle droit
    \draw[thick] (3.7,0) -- (3.7,0.3) -- (4,0.3);

    \node[below left] at (A) {$A$};
    \node[below right] at (B) {$B$};
    \node[above] at (S) {Sommet $S$};
    \node[below] at (H) {$H$};

    \draw[<->, >=stealth, slate700] (0,-0.4) -- (8,-0.4) node[midway, below] {Entrait $AB = 8\text{ m}$};
    \node at (2,-0.2) [below, font=\scriptsize] {$AH = 4\text{ m}$};
    \node at (6,-0.2) [below, font=\scriptsize] {$HB = 4\text{ m}$};
\end{tikzpicture}
\end{center}

Le triangle $ASB$ est isocèle en $S$. Le point $H$ est le milieu du segment $[AB]$ et le triangle $AHS$ est rectangle en $H$.

\begin{enumerate}
    \item \capSApproprier\ Justifier que la longueur $AH = 4\text{ m}$.
    \item \capRealiser\ Dans le triangle $AHS$ rectangle en $H$, calculer la longueur de l'arbalétrier $[AS]$. Arrondir au centième de mètre.
    \item \capRealiser\ Calculer l'angle de pente du toit $\widehat{SAH}$ formé avec l'horizontale. Arrondir au dixième de degré.
    \item \capValider\ Pour installer des tuiles en terre cuite, le DTU (Document Technique Unifié) exige une pente minimale de $30\degre$. La toiture est-elle conforme à cette réglementation ? Justifier.
\end{enumerate}
\lignesreponse[7]
\end{sujetdnb}

% ==============================================================================
% 7. CORRIGÉS DÉTAILLÉS
% ==============================================================================
\section{Corrigés Détaillés du Chapitre}

\begin{solution}[Corrigé des Exercices et du Sujet DNB]
\textbf{Exercice 1} :
Dans le triangle $RST$ rectangle en $S$, d'après le théorème de Pythagore :
$RT^2 = RS^2 + ST^2 = 4{,}8^2 + 6{,}4^2 = 23{,}04 + 40{,}96 = 64$.
Donc $RT = \sqrt{64} = \repbox{8\text{ cm}}$.

\vspace{2mm}
\textbf{Exercice 2} :
Le plus grand côté est $[KM]$ : $KM^2 = 25^2 = 625$. \\
D'autre part : $KL^2 + LM^2 = 7^2 + 24^2 = 49 + 576 = 625$. \\
On constate que $KM^2 = KL^2 + LM^2$. D'après la réciproque du théorème de Pythagore, le triangle $KLM$ est \textbf{rectangle en $L$}.

\vspace{2mm}
\textbf{Sujet DNB Pro (Charpente)} :
1. $H$ est le milieu de $[AB]$, donc $AH = 8 / 2 = \repbox{4\text{ m}}$. \\
2. Dans le triangle $AHS$ rectangle en $H$, d'après Pythagore : \\
$AS^2 = AH^2 + SH^2 = 4^2 + 2{,}4^2 = 16 + 5{,}76 = 21{,}76 \implies AS = \sqrt{21{,}76} \approx \repbox{4{,}66\text{ m}}$. \\
3. Dans le triangle $AHS$ rectangle en $H$ : $\tan(\widehat{SAH}) = \frac{SH}{AH} = \frac{2{,}4}{4} = 0{,}6$. \\
D'où $\widehat{SAH} = \arctan(0{,}6) \approx \repbox{31{,}0\degre}$. \\
4. L'angle calculé $31{,}0\degre \ge 30\degre$, la pente de la toiture est donc \textbf{pleinement conforme au DTU}.
\end{solution}


\chapter{Statistiques \& Probabilités}
\label{chap:stats_probas}

% ==============================================================================
% 1. OBJECTIFS & COMPÉTENCES DU RÉFÉRENTIEL
% ==============================================================================
\begin{cadreretenu}[Objectifs \& Compétences du DNB Pro]
\begin{itemize}
    \item \textbf{Statistiques} : Effectifs, fréquences, moyenne arithmétique simple et pondérée, médiane, étendue, diagrammes en bâtons et circulaires.
    \item \textbf{Probabilités} : Vocabulaire des événements (certain, impossible, contraire), calcul de probabilités simples, règle de la somme $P(A) + P(\bar{A}) = 1$, arbres de probabilités à deux épreuves.
    \item \textbf{Capacités} :
    \begin{itemize}
        \item \capRealiser\ Calculer la moyenne et la médiane d'une série statistique.
        \item \capAnalyser\ Comparer deux séries statistiques à l'aide des indicateurs de position et de dispersion.
        \item \capRealiser\ Déterminer la probabilité d'un événement sous forme de fraction irréductible, décimale ou pourcentage.
        \item \capValider\ Exploiter un arbre pondéré ou un tableau à double entrée pour dénombrer des issues.
    \end{itemize}
\end{itemize}
\end{cadreretenu}

\vspace{0.3cm}

% ==============================================================================
% 2. COURS : STATISTIQUES À UNE VARIABLE
% ==============================================================================
\section{Cours : Statistiques à Une Variable}

\subsection{Indicateurs de Position : Moyenne et Médiane}

\begin{cadredef}[Moyenne pondérée]
La \textbf{moyenne} pondérée d'une série statistique est le quotient de la somme de toutes les valeurs multipliées par leurs effectifs respectifs par l'effectif total $N$ :
\[ \bar{x} = \frac{n_1 x_1 + n_2 x_2 + \dots + n_p x_p}{N} = \frac{\sum n_i x_i}{N} \]
\end{cadredef}

\begin{cadredef}[Médiane et Étendue]
\begin{itemize}
    \item La \textbf{médiane} $Me$ est une valeur qui partage la série ordonnée en deux groupes de même effectif (au moins 50\,\% des valeurs sont inférieures ou égales à $Me$).
    \item L'\textbf{étendue} $E$ est la différence entre la plus grande valeur et la plus petite valeur :
    \[ E = x_{\max} - x_{\min} \]
\end{itemize}
\end{cadredef}

\begin{cadreexemple}[Calcul complet sur une série]
Contrôle qualité de 10 pièces usinées (longueur en mm) : $48\,;\; 49\,;\; 50\,;\; 50\,;\; 50\,;\; 51\,;\; 51\,;\; 52\,;\; 53\,;\; 56$.
\begin{itemize}
    \item \textbf{Moyenne} : $\bar{x} = \frac{48 + 49 + 3\times 50 + 2\times 51 + 52 + 53 + 56}{10} = \frac{510}{10} = \repbox{51\text{ mm}}$.
    \item \textbf{Médiane} : L'effectif total est $N = 10$ (pair). $10 / 2 = 5$. La médiane est la moyenne entre la 5\textsuperscript{e} valeur ($50$) et la 6\textsuperscript{e} valeur ($51$) : $Me = \frac{50+51}{2} = \repbox{50{,}5\text{ mm}}$.
    \item \textbf{Étendue} : $E = 56 - 48 = \repbox{8\text{ mm}}$.
\end{itemize}
\end{cadreexemple}

% ==============================================================================
% 3. COURS : PROBABILITÉS
% ==============================================================================
\section{Cours : Probabilités}

\begin{cadreprop}[Propriétés fondamentales]
Pour tout événement $A$ d'un univers $\Omega$ :
\begin{itemize}
    \item $0 \le P(A) \le 1$.
    \item En situation d'\textbf{équiprobabilité} :
    \[ P(A) = \frac{\text{Nombre d'issues favorables à } A}{\text{Nombre total d'issues possibles}} \]
    \item Événement contraire $\bar{A}$ : $P(\bar{A}) = 1 - P(A)$.
\end{itemize}
\end{cadreprop}

% ==============================================================================
% 4. QUESTIONS FLASH / AUTOMATISMES
% ==============================================================================
\section{Questions Flash / Automatismes}

\begin{automatisme}[Calculs Rapides (10 min)]
\noindent
\begin{tabularx}{\linewidth}{X|X}
\textbf{1.} Moyenne de $12$, $14$ et $16$ : $\bar{x} = \pointille[1.5cm]$ & \textbf{4.} On lance un dé équilibré à 6 faces. $P(\text{obtenir un nombre pair}) = \pointille[1.5cm]$ \\
\textbf{2.} Étendue de la série $\{ 3 ; 8 ; 15 ; 2 ; 19 \}$ : $E = \pointille[1.5cm]$ & \textbf{5.} Si $P(A) = 0{,}35$, alors $P(\bar{A}) = \pointille[1.5cm]$ \\
\textbf{3.} Médiane de la série ordonnée $\{ 4 ; 7 ; 9 ; 12 ; 15 \}$ : $Me = \pointille[1.5cm]$ & \textbf{6.} Une probabilité peut-elle être égale à $1{,}2$ ? \pointille[1.5cm] \\
\end{tabularx}
\end{automatisme}

% ==============================================================================
% 5. TRAVAUX DIRIGÉS (TD)
% ==============================================================================
\section{Travaux Dirigés (TD) : Exercices d'Entraînement}

\begin{exercice}[4 pts]{\capRealiser}{Contrôle Statistique en Logistique}
Un logisticien mesure les temps de préparation de 50 commandes dans un entrepôt (en minutes) :
\begin{center}
\begin{tabular}{|l|c|c|c|c|c|}
\hline
\rowcolor{slate200} \textbf{Temps (min)} & 5 & 10 & 15 & 20 & 25 \\
\hline
\textbf{Effectif} & 8 & 14 & 16 & 8 & 4 \\
\hline
\end{tabular}
\end{center}
\begin{enumerate}
    \item Calculer le temps moyen $\bar{t}$ de préparation d'une commande.
    \item Déterminer la classe médiane et l'étendue de cette série.
    \item Quel pourcentage de commandes est préparé en 15 minutes ou moins ?
\end{enumerate}
\lignesreponse[5]
\end{exercice}

\begin{exercice}[4 pts]{\capRealiser}{Tirage Aléatoire dans un Stock de Pièces}
Dans une boîte d'un atelier mécanique, on trouve :
\begin{itemize}
    \item 15 vis à tête fraisée (dont 5 défectueuses) ;
    \item 25 vis à tête cylindrique (dont 3 défectueuses).
\end{itemize}
On tire une vis au hasard dans la boîte.
\begin{enumerate}
    \item Quelle est la probabilité que la vis soit à tête fraisée ?
    \item Quelle est la probabilité que la vis soit défectueuse ?
    \item Quelle est la probabilité que la vis soit conforme (non défectueuse) ?
\end{enumerate}
\lignesreponse[5]
\end{exercice}

% ==============================================================================
% 6. SUJET DNB PRO
% ==============================================================================
\section{Vers le Brevet (DNB Pro)}

\begin{sujetdnb}[DNB Pro Métropole][9 pts]{Consommation Électrique \& Panneaux Photovoltaïques}
Une entreprise artisanale installe des panneaux solaires sur le toit de son atelier.
Le tableau suivant donne la production électrique mensuelle mesurée sur l'année (en kWh) :

\begin{center}
\begin{tabular}{|c|c|c|c|c|c|c|c|c|c|c|c|c|}
\hline
\rowcolor{slate200} \textbf{Mois} & J & F & M & A & M & J & J & A & S & O & N & D \\
\hline
\textbf{kWh} & 180 & 240 & 420 & 580 & 750 & 820 & 860 & 810 & 630 & 410 & 220 & 160 \\
\hline
\end{tabular}
\end{center}

\begin{enumerate}
    \item \capRealiser\ Calculer la production moyenne mensuelle d'électricité sur l'année.
    \item \capRealiser\ Déterminer la médiane et l'étendue de cette production.
    \item \capAnalyser\ On choisit un mois de l'année au hasard.
    \begin{enumerate}
        \item Quelle est la probabilité que la production dépasse $600\text{ kWh}$ ?
        \item Quelle est la probabilité que la production soit inférieure à $200\text{ kWh}$ ?
    \end{enumerate}
\end{enumerate}
\lignesreponse[6]
\end{sujetdnb}

% ==============================================================================
% 7. CORRIGÉS DÉTAILLÉS
% ==============================================================================
\section{Corrigés Détaillés du Chapitre}

\begin{solution}[Corrigé des Exercices et du Sujet DNB]
\textbf{Exercice 1 (Logistique)} :
1. Total effectif $N = 8 + 14 + 16 + 8 + 4 = 50$. \\
$\bar{t} = \frac{5\times 8 + 10\times 14 + 15\times 16 + 20\times 8 + 25\times 4}{50} = \frac{40 + 140 + 240 + 160 + 100}{50} = \frac{680}{50} = \repbox{13{,}6\text{ min}}$. \\
2. $50/2 = 25$. Effectifs cumulés : $8, 22, 38$. La médiane est de \textbf{15 min}. Étendue $E = 25 - 5 = \repbox{20\text{ min}}$. \\
3. Nombre de commandes $\le 15\text{ min}$ : $8 + 14 + 16 = 38$. Soit $\frac{38}{50} = \repbox{76\,\%}$.

\vspace{2mm}
\textbf{Exercice 2 (Stock de pièces)} :
Total de vis $= 15 + 25 = 40$. \\
1. $P(\text{fraisée}) = \frac{15}{40} = \frac{3}{8} = \repbox{0{,}375}$ (soit 37,5\,\%). \\
2. Total défectueuses $= 5 + 3 = 8$. $P(\text{défectueuse}) = \frac{8}{40} = \frac{1}{5} = \repbox{0{,}20}$ (soit 20\,\%). \\
3. $P(\text{conforme}) = 1 - 0{,}20 = \repbox{0{,}80}$ (soit 80\,\%).

\vspace{2mm}
\textbf{Sujet DNB Pro (Solaire)} :
1. Total annuel $= 180 + 240 + 420 + 580 + 750 + 820 + 860 + 810 + 630 + 410 + 220 + 160 = 6\,080\text{ kWh}$. \\
Moyenne mensuelle $= 6\,080 / 12 \approx \repbox{506{,}67\text{ kWh}}$. \\
2. Série ordonnée (12 valeurs) : $160, 180, 220, 240, 410, 420, 580, 630, 750, 810, 820, 860$. \\
Médiane : moyenne entre 6\textsuperscript{e} ($420$) et 7\textsuperscript{e} ($580$) : $Me = \frac{420+580}{2} = \repbox{500\text{ kWh}}$. \\
Étendue : $E = 860 - 160 = \repbox{700\text{ kWh}}$. \\
3. (a) Mois $> 600\text{ kWh}$ : Juin, Juillet, Août, Septembre, Mai (soit 5 mois). $P = \frac{5}{12} \approx \repbox{0{,}417}$. \\
(b) Mois $< 200\text{ kWh}$ : Janvier, Décembre (soit 2 mois). $P = \frac{2}{12} = \frac{1}{6} \approx \repbox{0{,}167}$.
\end{solution}


\chapter{Aires, Volumes \& Conversions d'Unités}
\label{chap:aires_volumes}

% ==============================================================================
% 1. OBJECTIFS & COMPÉTENCES DU RÉFÉRENTIEL
% ==============================================================================
\begin{cadreretenu}[Objectifs \& Compétences du DNB Pro]
\begin{itemize}
    \item \textbf{Connaissances} : Formules d'aires usuelles (rectangle, triangle, disque, trapèze), formules de volumes (pavé droit, prisme droit, cylindre, cône, pyramide, boule), équivalence $1\text{ L} = 1\text{ dm}^3$.
    \item \textbf{Capacités} :
    \begin{itemize}
        \item \capRealiser\ Effectuer des conversions d'unités de longueurs, d'aires et de volumes/capacités.
        \item \capRealiser\ Calculer l'aire d'une surface complexe en la décomposant en figures simples.
        \item \capRealiser\ Calculer le volume d'un réservoir, d'un silo ou d'une cuve industrielle.
        \item \capAnalyser\ Déterminer une masse à partir d'un volume et d'une masse volumique ($\rho = \frac{m}{V}$).
    \end{itemize}
\end{itemize}
\end{cadreretenu}

\vspace{0.3cm}

% ==============================================================================
% 2. COURS : FORMULAIRE OFFICIEL AIRES ET VOLUMES
% ==============================================================================
\section{Cours : Formulaire d'Aires et de Volumes}

\subsection{Surfaces Planes Usuelles}

\noindent
\begin{tabularx}{\linewidth}{|X|X|X|}
\hline
\rowcolor{colDefBg} \textbf{Rectangle} & \textbf{Triangle} & \textbf{Disque} \\
\hline
\begin{center}
\begin{tikzpicture}[scale=0.6]
    \draw[fill=colDef!20, draw=colDef, thick] (0,0) rectangle (3,1.8);
    \node at (1.5,0.9) {$\mathcal{A} = L \times \ell$};
\end{tikzpicture}
\end{center} &
\begin{center}
\begin{tikzpicture}[scale=0.6]
    \draw[fill=colProp!20, draw=colProp, thick] (0,0) -- (3,0) -- (1,2) -- cycle;
    \draw[dashed] (1,2) -- (1,0);
    \node at (1.5,0.7) {$\mathcal{A} = \frac{b \times h}{2}$};
\end{tikzpicture}
\end{center} &
\begin{center}
\begin{tikzpicture}[scale=0.6]
    \draw[fill=colExem!20, draw=colExem, thick] (0,0) circle (1.2);
    \draw[->] (0,0) -- (1.2,0) node[midway, above, font=\tiny] {$R$};
    \node at (0,-0.3) [font=\footnotesize] {$\mathcal{A} = \pi R^2$};
\end{tikzpicture}
\end{center} \\
\hline
\end{tabularx}

\subsection{Solides et Volumes Usuels}

\begin{cadreprop}[Volumes des solides de l'espace]
\begin{itemize}
    \item \textbf{Pavé droit (Parallélépipède)} : $\mathcal{V} = L \times \ell \times h$
    \item \textbf{Prisme droit \& Cylindre de révolution} : $\mathcal{V} = \mathcal{A}_{\text{base}} \times h = \pi R^2 h$
    \item \textbf{Pyramide \& Cône de révolution} : $\mathcal{V} = \dfrac{1}{3} \times \mathcal{A}_{\text{base}} \times h = \dfrac{1}{3} \pi R^2 h$
    \item \textbf{Boule / Sphère} : $\mathcal{V} = \dfrac{4}{3} \pi R^3$ \quad et \quad $\mathcal{A}_{\text{sphère}} = 4 \pi R^2$
\end{itemize}
\end{cadreprop}

\begin{cadreretenu}[Conversions d'Unités Incontournables]
\[ 1\text{ m}^3 = 1\,000\text{ dm}^3 = 1\,000\text{ L} \qquad\text{et}\qquad 1\text{ dm}^3 = 1\text{ L} = 1\,000\text{ cm}^3 = 1\,000\text{ mL} \]
\end{cadreretenu}

% ==============================================================================
% 3. QUESTIONS FLASH / AUTOMATISMES
% ==============================================================================
\section{Questions Flash / Automatismes}

\begin{automatisme}[Conversions \& Formules (10 min)]
\noindent
\begin{tabularx}{\linewidth}{X|X}
\textbf{1.} $2{,}5\text{ m}^3 = \pointille[1.5cm]\text{ L}$ & \textbf{4.} Volume d'un cube d'arête $a = 4\text{ cm}$ : $V = \pointille[1.5cm]\text{ cm}^3$ \\
\textbf{2.} $750\text{ cm}^3 = \pointille[1.5cm]\text{ L}$ & \textbf{5.} Aire d'un disque de rayon $R = 10\text{ cm}$ : $\mathcal{A} \approx \pointille[1.5cm]\text{ cm}^2$ \\
\textbf{3.} $12\text{ m}^2 = \pointille[1.5cm]\text{ dm}^2$ & \textbf{6.} Masse volumique de l'eau : $1\text{ L d'eau} = \pointille[1cm]\text{ kg}$ \\
\end{tabularx}
\end{automatisme}

% ==============================================================================
% 4. TRAVAUX DIRIGÉS (TD)
% ==============================================================================
\section{Travaux Dirigés (TD) : Exercices d'Entraînement}

\begin{exercice}[4 pts]{\capRealiser}{Cuve de Récupération d'Eau de Pluie}
Une cuve cylindrique de récupération d'eau a un diamètre intérieur de $D = 1{,}20\text{ m}$ et une hauteur utile de $h = 1{,}80\text{ m}$.
\begin{enumerate}
    \item Calculer le rayon intérieur $R$ en mètres.
    \item Calculer le volume de la cuve en $\text{m}^3$. Arrondir au centième.
    \item En déduire la capacité maximale de la cuve en Litres.
\end{enumerate}
\lignesreponse[4]
\end{exercice}

\begin{exercice}[4 pts]{\capAnalyser}{Coulage d'une Dalle en Béton (BTP)}
Un maçon réalise une dalle en béton rectangulaire de longueur $L = 8\text{ m}$, de largeur $\ell = 4{,}50\text{ m}$ et d'épaisseur $e = 15\text{ cm}$.
\begin{enumerate}
    \item Exprimer l'épaisseur $e$ en mètres.
    \item Calculer le volume de béton nécessaire en $\text{m}^3$.
    \item Sachant que la masse volumique du béton armé est $\rho = 2\,400\text{ kg/m}^3$, calculer la masse totale de la dalle en tonnes.
\end{enumerate}
\lignesreponse[5]
\end{exercice}

% ==============================================================================
% 5. VERS LE BREVET (DNB PRO)
% ==============================================================================
\section{Vers le Brevet (DNB Pro)}

\begin{sujetdnb}[DNB Pro Métropole][9 pts]{Traitement d'une Piscine Municipale}
Une piscine municipale dispose d'un bassin de forme pavé droit de longueur $25\text{ m}$, de largeur $10\text{ m}$ et de profondeur constante $1{,}60\text{ m}$.

\begin{enumerate}
    \item \capRealiser\ Calculer le volume d'eau contenu dans ce bassin en $\text{m}^3$, puis en Litres.
    \item \capAnalyser\ Pour traiter l'eau, le maître-nageur doit verser $2\text{ cL}$ de produit désinfectant pour $1\,000\text{ L}$ d'eau.
    Calculer la quantité totale de désinfectant à verser (en Litres).
    \item \capValider\ La pompe de filtration a un débit de $40\text{ m}^3/\text{h}$. Combien d'heures faut-il pour renouveler la totalité de l'eau du bassin ?
\end{enumerate}
\lignesreponse[6]
\end{sujetdnb}

% ==============================================================================
% 6. CORRIGÉS DÉTAILLÉS
% ==============================================================================
\section{Corrigés Détaillés du Chapitre}

\begin{solution}[Corrigé des Exercices et du Sujet DNB]
\textbf{Exercice 1 (Cuve)} :
1. $R = 1{,}20 / 2 = \repbox{0{,}60\text{ m}}$. \\
2. $V = \pi \times R^2 \times h = \pi \times 0{,}60^2 \times 1{,}80 \approx \repbox{2{,}04\text{ m}^3}$. \\
3. $1\text{ m}^3 = 1\,000\text{ L} \implies 2{,}04 \times 1\,000 = \repbox{2\,040\text{ Litres}}$.

\vspace{2mm}
\textbf{Exercice 2 (Béton BTP)} :
1. $e = 15\text{ cm} = \repbox{0{,}15\text{ m}}$. \\
2. $V = 8 \times 4{,}50 \times 0{,}15 = \repbox{5{,}40\text{ m}^3}$. \\
3. Masse $= 5{,}40 \times 2\,400 = 12\,960\text{ kg} = \repbox{12{,}96\text{ tonnes}}$.

\vspace{2mm}
\textbf{Sujet DNB Pro (Piscine)} :
1. $V = 25 \times 10 \times 1{,}60 = \repbox{400\text{ m}^3} = \repbox{400\,000\text{ Litres}}$. \\
2. Produit $= \frac{400\,000}{1\,000} \times 2\text{ cL} = 400 \times 2 = 800\text{ cL} = \repbox{8\text{ Litres}}$. \\
3. Temps de filtration $= \frac{400}{40} = \repbox{10\text{ heures}}$.
\end{solution}


% ==============================================================================
% QUATRIÈME DE COUVERTURE (PAGE DE FIN DU LIVRET)
% ==============================================================================
\newpage
\thispagestyle{empty}
\begin{tikzpicture}[remember picture, overlay]
    \fill[slate900] (current page.north west) rectangle (current page.south east);
    
    \node[anchor=north, text=colSecondary, font=\sffamily\bfseries\Large] at ([yshift=-3cm]current page.north) {
        FASCICULE PÉDAGOGIQUE MATHÉMATIQUES 3PM
    };
    
    \node[anchor=center, text width=15cm, align=center, text=slate200, font=\sffamily\normalsize] at (current page.center) {
        \begin{tcolorbox}[
            enhanced,
            colback=slate800,
            colframe=colPrimary,
            arc=8pt,
            boxrule=1.5pt,
            coltext=white,
            title={\faCheckCircle\ \textbf{MÉMO DU RÉFÉRENTIEL DNB PRO}},
            fonttitle=\sffamily\bfseries\large\color{colSecondary}
        ]
        \small
        \begin{itemize}
            \item \textbf{Calcul Littéral \& Équations} : Distributivité simple et double, identité $a^2-b^2$, équations $ax+b=cx+d$ et produit nul.
            \item \textbf{Fonctions} : Fonctions linéaires (droite par l'origine) et affines $f(x)=ax+b$.
            \item \textbf{Géométrie} : Théorème de Pythagore ($BC^2 = AB^2 + AC^2$), Trigonométrie (\textbf{CAH-SOH-TOA}).
            \item \textbf{Statistiques \& Probabilités} : Moyenne, médiane, étendue, équiprobabilité ($P(A) = \frac{\text{favorables}}{\text{totales}}$).
            \item \textbf{Grandeurs \& Mesures} : $1\text{ L} = 1\text{ dm}^3$, $1\text{ m}^3 = 1\,000\text{ L}$, volumes prismes, cylindres, cônes et boules.
        \end{itemize}
        \end{tcolorbox}
        
        \vspace{1.5cm}
        {\large\bfseries Conforme aux Programmes Officiels de l'Éducation Nationale}\\[2mm]
        {\small Réalisé pour une utilisation autonome ou collective en classe et sur Overleaf.}
    };
    
    \node[anchor=south, text=slate500, font=\sffamily\footnotesize] at ([yshift=2cm]current page.south) {
        \thelycee\ -- \theannee\ $\cdot$ Tous droits réservés.
    };
\end{tikzpicture}

\end{document}
